% ============================================================
%  第4章：对偶单纯形法与灵敏度分析
% ============================================================
\section{第4章 \quad 对偶单纯形法与灵敏度分析}

\subsection{4.1 为什么需要对偶单纯形法}

第2章的（原始）单纯形法走的是"原始可行 $\to$ 对偶可行"的路线：始终保持 $b \ge 0$（原始可行），通过迭代逐步让检验数全部 $\le 0$（达到对偶可行，即最优）。

但在实际中，我们经常遇到\emph{反过来}的情况：已经有一张检验数全部 $\le 0$ 的表（对偶可行），但某些 $b_i < 0$（原始不可行）。此时重头用大M法太浪费——对偶单纯形法（Dual Simplex Method）正是为这种场景设计的。

\begin{defbox}
\textbf{对偶单纯形法的核心思想}

保持\emph{对偶可行性}（检验数 $\le 0$），逐步消除 $b$ 列的负值，最终达到原始可行——此时既是原始可行又是对偶可行，即为最优解。

\textbf{典型适用场景}：
\begin{itemize}
  \item 灵敏度分析中 $b$ 改变导致基变量取值变负
  \item 新增一个约束后原最优解不再满足
  \item 初始表中检验数已 $\le 0$，但某些 $b_i < 0$（例如 $\ge$ 约束标准化时乘以 $-1$）
\end{itemize}
\end{defbox}

\medskip
PPT第6.19页阐述了对偶单纯形法的基本思想：与原始单纯形法"保原始可行、追对偶可行"对称，对偶单纯形法"保对偶可行、追原始可行"，两者构成完美的镜像。
\refslide{6对偶理论}{25}{6.19}
\medskip

\subsection{4.2 对偶单纯形法操作步骤}

\begin{defbox}
\textbf{方法 4.1（对偶单纯形法一次迭代）}

前提：表中所有检验数 $\le 0$（对偶可行），但存在 $\bar{b}_r < 0$（原始不可行）。

\begin{enumerate}
  \item \textbf{离基选择}：选择 $b$ 列中\emph{最负}值所在行的基变量离基，记为第 $r$ 行。
  \item \textbf{进基选择（对偶最小比值）}：在第 $r$ 行中，对所有\emph{负系数} $\bar{a}_{rj} < 0$，计算
        $\displaystyle \theta_j = \frac{\sigma_j}{|\bar{a}_{rj}|} = \frac{\sigma_j}{-\bar{a}_{rj}}$。
        取 $\theta_j$ \emph{最小}的变量 $x_k$ 进基。\\
        若该行全无负系数，则原问题\emph{无可行解}。
  \item \textbf{转轴}：以 $\bar{a}_{rk}$ 为主元进行行变换，使 $x_k$ 列成为单位向量，同时更新检验数行。
\end{enumerate}
\end{defbox}

\medskip
PPT第6.20页给出了对偶单纯形法的详细推导，包括离基变量和进基变量的选取准则、主元消去后对偶可行性的保持证明。
\refslide{6对偶理论}{26}{6.20}
\medskip

\begin{notebox}
\textbf{原始单纯形 vs 对偶单纯形对比}：
\medskip
{\small
\begin{tabular}{@{}lcc@{}}
\toprule
& \textbf{原始单纯形} & \textbf{对偶单纯形} \\
\midrule
先选什么 & 进基（最大正检验数） & 离基（最负 $b$ 值） \\
后选什么 & 离基（$\min b_i/a_{ik}$） & 进基（$\min |\sigma_j|/|a_{rj}|$） \\
比值在哪儿算 & 进基列（纵向） & 离基行（横向） \\
保持什么 & 原始可行（$b\ge 0$） & 对偶可行（$\sigma \le 0$） \\
推动什么 & 检验数全部 $\le 0$ & $b$ 列全部 $\ge 0$ \\
\bottomrule
\end{tabular}
}
\end{notebox}

\subsection{4.3 为什么是这个顺序——因果解释}

\textbf{为什么先选离基（最负 $b$）？}

对偶单纯形法的目标是消除 $b$ 列的负值。最负的 $b_r$ 意味着第 $r$ 个约束被违反得最严重——这个基变量取值为负，离可行域最远。优先修正它能最快地向可行域靠近。这和对偶问题中目标函数是 $\max\ b^T w$ 有关：$b_r < 0$ 时让 $x_{B_r}$ 离基，等价于对偶问题中让目标函数值增加——这正是对偶问题要做的（最大化）。

\textbf{为什么在负行系数上取最小比值？}

假设第 $r$ 行上 $\bar{a}_{rj} < 0$。当 $x_j$ 进基时，第 $r$ 个基变量的值变为：
\[
\bar{b}_r' = \frac{\bar{b}_r}{\bar{a}_{rj}} \quad (\text{因为主元消去后新 } b_r \text{ 为原值除以主元})
\]
更精确地说，转轴后 $b_r$ 变为 $\bar{b}_r / \bar{a}_{rj} > 0$（因 $\bar{b}_r < 0$ 且 $\bar{a}_{rj} < 0$）。但其他行 $i \neq r$ 的 $b_i$ 也会受影响：$\bar{b}_i' = \bar{b}_i - \frac{\bar{a}_{ij}}{\bar{a}_{rj}}\bar{b}_r$。取 $\theta_j = |\sigma_j|/|\bar{a}_{rj}|$ 最小，相当于在对偶问题中做"最保守的步子"——保证所有检验数都保持在 $\le 0$，不破坏对偶可行性。这与原始单纯形法中"取最小 $b_i/a_{ik}$ 以保证不跑出可行域"的几何直觉完全对称。

\textbf{两者的对称性总结}：原始单纯形在对偶不可行的方向上试探，用最小比值保护原始可行性；对偶单纯形在原始不可行的方向上试探，用最小比值保护对偶可行性。

\subsection{4.4 自编例子：对偶单纯形法完整迭代}

下面用一个自编的简单 LP 展示对偶单纯形法的完整过程，并与原始单纯形法的迭代思路对照。

\textbf{问题}：
\[
\begin{aligned}
\min\ & x_1 + 2x_2 \\
\text{s.t.}\ & x_1 + x_2 \ge 3 \\
& 2x_1 + x_2 \ge 4 \\
& x_1, x_2 \ge 0
\end{aligned}
\]

\textbf{标准化}（引入剩余变量 $x_3,x_4 \ge 0$）：
\[
\begin{aligned}
\min\ & x_1 + 2x_2 \\
\text{s.t.}\ & x_1 + x_2 - x_3 = 3 \\
& 2x_1 + x_2 - x_4 = 4 \\
& x_j \ge 0,\ j=1,\dots,4
\end{aligned}
\]

两个约束都是 $\ge$ 型，剩余变量系数为 $-1$，没有现成的初始基。将约束两边乘 $-1$，得：
\[
\begin{aligned}
-x_1 - x_2 + x_3 &= -3 \\
-2x_1 - x_2 + x_4 &= -4
\end{aligned}
\]
此时 $x_3,x_4$ 构成初始基，但 $b = (-3,-4)^T$ 为负。检验数呢？$c_B = (0,0)$，全部 $\sigma_j = 0 - c_j \le 0$：$\sigma_1 = -1$，$\sigma_2 = -2$，$\sigma_3 = \sigma_4 = 0$。这是一张典型的\emph{对偶可行但原始不可行}的表——对偶单纯形法的起点。

\medskip
\textbf{初始表（表 4.1）}：
\begingroup
\footnotesize
\setlength{\tabcolsep}{4pt}
\[
\begin{array}{c|cccc|c}
\text{基} & x_1 & x_2 & x_3 & x_4 & b \\ \toprule
x_3 & -1 & -1 & 1 & 0 & -3 \\
x_4 & -2 & -1 & 0 & 1 & -4 \\ \midrule
\sigma_j & -1 & -2 & 0 & 0 & \text{---}
\end{array}
\]
\endgroup

\medskip
\textbf{迭代 1}：
\begin{itemize}
  \item \textbf{离基}：$b$ 列最负为 $-4$（第 2 行），$x_4$ 离基。
  \item \textbf{进基}：第 2 行负系数有 $x_1(-2)$、$x_2(-1)$。\\
        $\theta_1 = |-1|/|-2| = 0.5$，$\theta_2 = |-2|/|-1| = 2$。$\min = 0.5$，$x_1$ 进基。主元 $=-2$。
  \item \textbf{转轴}：第 2 行 $\div (-2)$，再用它消去其他行的 $x_1$ 列。
\end{itemize}

\begingroup
\footnotesize
\setlength{\tabcolsep}{4pt}
\[
\begin{array}{c|cccc|c}
\text{基} & x_1 & x_2 & x_3 & x_4 & b \\ \toprule
x_3 & 0 & -1/2 & 1 & -1/2 & -1 \\
x_1 & 1 & 1/2  & 0 & -1/2 & 2 \\ \midrule
\sigma_j & 0 & -3/2 & 0 & -1/2 & \text{---}
\end{array}
\]
\endgroup

所有 $\sigma \le 0$（对偶可行保持），但 $b_1 = -1$ 仍为负。

\medskip
\textbf{迭代 2}：
\begin{itemize}
  \item \textbf{离基}：$b_1 = -1$，$x_3$ 离基。
  \item \textbf{进基}：第 1 行负系数 $x_2(-1/2)$、$x_4(-1/2)$。\\
        $\theta_2 = (3/2)/(1/2) = 3$，$\theta_4 = (1/2)/(1/2) = 1$。$\min = 1$，$x_4$ 进基。主元 $=-1/2$。
  \item \textbf{转轴}：第 1 行 $\times (-2)$ 后消去。
\end{itemize}

\begingroup
\footnotesize
\setlength{\tabcolsep}{4pt}
\[
\begin{array}{c|cccc|c}
\text{基} & x_1 & x_2 & x_3 & x_4 & b \\ \toprule
x_4 & 0 & 1 & -2 & 1 & 2 \\
x_1 & 1 & 1 & -1 & 0 & 3 \\ \midrule
\sigma_j & 0 & -1 & -1 & 0 & \text{---}
\end{array}
\]
\endgroup

$b \ge 0$ 且 $\sigma \le 0$，达到最优！最优解 $x^* = (3, 0, 0, 2)^T$，$z^* = 1 \times 3 + 2 \times 0 = 3$。

\begin{notebox}
\textbf{对称性观察}：如果这个 LP 变个方向用原始单纯形法，需要先找初始 BFS（比如用大M法加人工变量），走"保 $b\ge 0$ 追 $\sigma\le 0$"的路线。而对偶单纯形法\emph{利用了对偶可行性}，省掉了人工变量——因为 $\ge$ 约束乘 $-1$ 后自动给出对偶可行的初始表。这正是对偶单纯形法的实用价值。
\end{notebox}

\subsection{4.5 灵敏度分析}

灵敏度分析（Sensitivity Analysis）回答一个实际问题：当 LP 的参数（价格 $c$、资源量 $b$、工艺系数 $A$）发生变化时，最优解会怎么变？我们从\emph{已求出的最优表}出发分析，而非重头求解。

\begin{defbox}
\textbf{灵敏度分析总决策表}
\medskip
{\footnotesize
\setlength{\tabcolsep}{2.5pt}
\begin{tabular}{@{}p{36mm}p{34mm}p{46mm}@{}}
\toprule
原问题状态 & 对偶问题状态 & 应对措施 \\
\midrule
可行 + 对偶可行 & 对偶可行 & 最优解/最优基不变 \\
可行 + 对偶不可行 & 对偶不可行 & 原始单纯形继续迭代 \\
不可行 + 对偶可行 & 对偶可行 & 对偶单纯形继续迭代 \\
不可行 + 对偶不可行 & 对偶不可行 & 加人工变量从头求解 \\
\bottomrule
\end{tabular}
}
\end{defbox}

\subsubsection{4.5.1 目标系数 $c$ 的变化}

PPT第6.46页讲解了改变系数向量 $c$ 时的灵敏度分析框架：分为非基变量和基变量两种情况。PPT第6.47页进一步说明 $c$ 变化时只需计算新的检验数，其余不变。
\refslide{6对偶理论}{58}{6.46}
\refslide{6对偶理论}{59}{6.47}
\medskip

\textbf{情况一：非基变量 $x_j$ 的 $c_j$ 变为 $c_j'$。}

只有 $x_j$ 自己的检验数受影响：
\[
\sigma_j' = z_j - c_j' = \sigma_j - (c_j' - c_j)
\]
若 $\sigma_j' \le 0$，最优解不变；若 $\sigma_j' > 0$，以当前表为起点，$x_j$ 进基继续单纯形迭代。

\begin{notebox}
\textbf{为什么只影响一个检验数？} 检验数的计算公式是 $\sigma_j = c_B^T B^{-1} a_j - c_j$。非基变量的 $c_j$ 只出现在公式的第 2 项。$c_B$（基变量价格向量）不变，$B^{-1}a_j$ 不变，所以 $z_j = c_B^T B^{-1}a_j$ 不变——唯一变化的是减去新 $c_j'$。而基变量 $c$ 变化时 $c_B$ 改变，所有 $z_j$ 都要重新算。
\end{notebox}

\textbf{情况二：基变量 $x_k$ 的 $c_k$ 变为 $c_k'$。}

$c_B$ 变了，\emph{所有}非基变量的检验数须全部重新计算：
\[
\sigma_j' = (c_B')^T B^{-1} a_j - c_j, \quad \forall j \in \text{非基}
\]
还须更新目标值：$z_0' = (c_B')^T B^{-1} b$。若所有 $\sigma_j' \le 0$，仍最优；否则继续迭代。

\textbf{例：} 在 \S4.4 的最优表中，$c_1$ 由 1 变为 3。$x_1$ 是基变量（$c_B$ 中 $c_1$ 变了），重算：
$c_B' = (c_4=0, c_1=3)$，$\sigma_2' = (0,3) \cdot (1,-1)^T - 2 = -3 - 2 = -5 \le 0$，
$\sigma_3' = (0,3) \cdot (-2,-1)^T - 0 = -3 \le 0$。仍最优。但若 $c_1$ 由 1 变为 0：
$\sigma_2' = (0,0) \cdot (1,-1)^T - 2 = -2 \le 0$，$\sigma_3' = (0,0) \cdot (-2,-1)^T = 0 \le 0$，仍最优。

\subsubsection{4.5.2 右端项 $b$ 的变化}

$b$ 变为 $b'$ 后，新右端 $\bar{b}' = B^{-1} b'$。
\begin{itemize}
  \item 若 $\bar{b}' \ge 0$：最优基不变，只是基变量取值变了。新目标值 $z' = c_B^T \bar{b}'$。
  \item 若出现 $\bar{b}_i' < 0$：检验数仍 $\le 0$（$b$ 变化不影响 $\sigma$），但对偶可行而原始不可行 —— 用对偶单纯形法恢复可行性。
\end{itemize}

\subsubsection{4.5.3 新增约束}

PPT第6.57页讲解加入新约束的处理流程：先检验原最优解是否满足新约束；若不满足，将新约束添入最优表，消去基变量在该行的影响后，用对偶单纯形法继续求解。
\refslide{6对偶理论}{71}{6.57}
\medskip

步骤：
\begin{enumerate}
  \item 将原最优解代入新约束。若满足 → 最优解不变，结束。
  \item 若不满足：引入松弛/剩余变量，将新约束作为新行添入最优表。
  \item 在新行中消去基变量（用现有的约束行消去基变量在新行中的系数）。
  \item 此时检验数仍 $\le 0$（新松弛变量检验数为 0），但新行 $b$ 可能为负 → 用对偶单纯形法。
\end{enumerate}

\subsubsection{4.5.4 新增决策变量}

新变量 $x_{n+1}$ 的 $c_{n+1}$ 和列 $a_{n+1}$ 已知。计算 $B^{-1}a_{n+1}$ 代入表，计算 $\sigma_{n+1} = c_B^T B^{-1} a_{n+1} - c_{n+1}$。若 $\le 0$，原解仍最优；否则新变量进基继续迭代。

\subsection{4.6 自编例题：灵敏度分析全流程}

\textbf{问题（自编）}：某工厂生产甲、乙两种产品，LP 模型如下：
\[
\begin{aligned}
\min\ & -2x_1 - 3x_2 \quad (\text{即 }\max\ 2x_1+3x_2) \\
\text{s.t.}\ & x_1 + 2x_2 \le 8 \quad \text{（原料A）} \\
& 3x_1 + x_2 \le 9 \quad \text{（原料B）} \\
& x_1, x_2 \ge 0
\end{aligned}
\]
引入松弛变量 $x_3,x_4$ 化为标准型后，用单纯形法求得\emph{最优表 4.2}：

\begingroup
\footnotesize
\setlength{\tabcolsep}{4pt}
\[
\begin{array}{c|cccc|c}
\text{基} & x_1 & x_2 & x_3   & x_4   & b \\ \toprule
x_2 & 0 & 1 & 3/5    & -1/5  & 3 \\
x_1 & 1 & 0 & -1/5   & 2/5   & 2 \\ \midrule
\sigma_j & 0 & 0 & -7/5 & -1/5 & \text{---}
\end{array}
\]
\endgroup

最优解 $x^* = (2,3,0,0)^T$，$z^* = -13$（利润 $=13$）。$B^{-1} = \begin{bmatrix} 3/5 & -1/5 \\ -1/5 & 2/5 \end{bmatrix}$（从 $x_3,x_4$ 列读得）。

\medskip
\textbf{情景A：非基变量 $c_3$ 变化。}
$x_3$ 是松弛变量，原 $c_3 = 0$。若原料A有了成本，$c_3$ 变为 $1$：
\[
\sigma_3' = \sigma_3 - (c_3' - c_3) = -\frac{7}{5} - (1 - 0) = -\frac{12}{5} \le 0
\]
最优解不变。\emph{解释：$x_3$ 是松弛变量，当前=0，提高其成本只会让它更不值得入基，最优解不受影响。读者可验证：只有当 $c_3' - c_3 \ge |\sigma_3| = 7/5$，即 $c_3$ 提高到 $7/5$ 以上时，$\sigma_3'$ 才变正，$x_3$ 才应进基。}

\medskip
\textbf{情景B：基变量 $c_1$ 变化。}
产品甲的价格变化，$c_1$ 从 $-2$ 变为 $-1$（利润从 2 降到 1）。$x_1$ 是基变量，$c_B$ 改变，重算所有非基检验数：
\[
c_B' = (c_2=-3,\ c_1=-1)
\]
\[
\sigma_3' = (-3,-1) \cdot \begin{bmatrix}3/5\\-1/5\end{bmatrix} - 0 = -\frac{9}{5} + \frac{1}{5} = -\frac{8}{5} \le 0
\]
\[
\sigma_4' = (-3,-1) \cdot \begin{bmatrix}-1/5\\2/5\end{bmatrix} - 0 = \frac{3}{5} - \frac{2}{5} = \frac{1}{5} > 0
\]
$\sigma_4' > 0$，最优解改变！需以当前表为起点，$x_4$ 进基继续迭代（略）。

\medskip
\textbf{情景C：右端 $b$ 变化（基不变）。}
原料A从 8 增到 10（$b_1' = 10$）：
\[
\bar{b}' = B^{-1} b' = \begin{bmatrix}3/5 & -1/5\\-1/5 & 2/5\end{bmatrix}
\begin{bmatrix}10\\9\end{bmatrix}
= \begin{bmatrix}21/5\\8/5\end{bmatrix} \ge 0
\]
$b' \ge 0$，最优基不变。新解：$x_2=21/5,\ x_1=8/5$，$z' = -79/5$（利润 $15.8$）。

\medskip
\textbf{情景D：右端 $b$ 变化（需对偶单纯形）。}
原料A从 8 减到 2（$b_1' = 2$）：
\[
\bar{b}' = B^{-1} b' = \begin{bmatrix}3/5 & -1/5\\-1/5 & 2/5\end{bmatrix}
\begin{bmatrix}2\\9\end{bmatrix}
= \begin{bmatrix}-3/5\\16/5\end{bmatrix}
\]
$\bar{b}_1 = -3/5 < 0$！原基不可行，但检验数仍 $\le 0$——恰好用对偶单纯形法：

\begingroup
\footnotesize
\setlength{\tabcolsep}{4pt}
\[
\begin{array}{c|cccc|c}
\text{基} & x_1 & x_2 & x_3 & x_4 & b \\ \toprule
x_2 & 0 & 1 & 3/5  & -1/5 & -3/5 \\
x_1 & 1 & 0 & -1/5 & 2/5  & 16/5 \\ \midrule
\sigma_j & 0 & 0 & -7/5 & -1/5 & \text{---}
\end{array}
\]
\endgroup

$x_2$ 离基（$b_1 = -3/5$ 唯一负）。第 1 行负系数仅 $x_4(-1/5)$，$\theta_4 = (1/5)/(1/5) = 1$，$x_4$ 进基。转轴后：
\begingroup
\footnotesize
\setlength{\tabcolsep}{4pt}
\[
\begin{array}{c|cccc|c}
\text{基} & x_1 & x_2 & x_3 & x_4 & b \\ \toprule
x_4 & 0 & -5 & -3 & 1 & 3 \\
x_1 & 1 & 2  & 1  & 0 & 2 \\ \midrule
\sigma_j & 0 & -1 & -2 & 0 & \text{---}
\end{array}
\]
\endgroup
新最优解 $x^* = (2,0,0,3)^T$，$z^* = -4$（利润 $4$）。原料A锐减后，产品乙停产，只剩产品甲 $x_1=2$。

\subsection{4.7 真题：第一题(3)——非基变量 $c$ 变化}

\begin{notebox}
\textbf{【2025 真题·第一题(3)】} 第一题求得最优解 $x^* = (0,\,11/3,\,7/3)^T$（见第2章表 2.3）后，若 $c_1$ 由 2 改为 0，原最优解是否仍为最优解？
\end{notebox}

\medskip

在表 2.3 中，$x_1$ 是\emph{非基变量}（取值为 $0$），其当前检验数 $\sigma_1 = -7/3$（见第2章最优表的末行首列）。这是典型的"情况一"——非基变量的 $c$ 变化只影响自己的检验数。

当 $c_1$ 从 2 变为 0 时：
\[
\sigma_1' = \sigma_1 - (c_1' - c_1) = -\frac{7}{3} - (0 - 2) = -\frac{7}{3} + 2 = -\frac{1}{3} \le 0
\]
检验数仍 $\le 0$，故原最优解不变。

\emph{注：此题体现了灵敏度分析的考试典型套路——不需要重跑单纯形法，只需从最优表中提取非基变量的检验数，用一步公式即可判断。}

\subsection{4.8 真题：第二题——新增约束+对偶单纯形法}

\begin{notebox}
\textbf{【2025 真题·第二题】} 已知某 LP 的初始和最优单纯形表如下。加入新约束 $2x_1 - x_2 + x_3 \le 9$。原最优解是否满足新约束？若不满足，求新最优解。
\end{notebox}

\medskip

\textbf{题目给定的初始表}：
\begingroup
\footnotesize
\setlength{\tabcolsep}{3pt}
\[
\begin{array}{c|ccccc|c}
\text{基} & x_1 & x_2 & x_3 & x_4 & x_5 & b \\ \toprule
x_4 & 2 & 3 & 1 & 1 & 0 & 4 \\
x_5 & 5 & 4 & 3 & 0 & 1 & 11 \\ \midrule
z_j-c_j & 9 & 8 & 5 & 0 & 0 & 0
\end{array}
\]
\endgroup

\textbf{题目给定的最优表}：
\begingroup
\footnotesize
\setlength{\tabcolsep}{3pt}
\[
\begin{array}{c|ccccc|c}
\text{基} & x_1 & x_2 & x_3 & x_4 & x_5 & b \\ \toprule
x_1 & 1 & 5 & 0 & 3 & -1 & 1 \\
x_3 & 0 & -7 & 1 & -5 & 2 & 2 \\ \midrule
z_j-c_j & 0 & -2 & 0 & -2 & -1 & -19
\end{array}
\]
\endgroup

原最优解：$x^* = (1,\,0,\,2,\,0,\,0)^T$，$z^* = 19$（注意：题目检验数行末是 $-19$，故 $z^* = 19$）。

\medskip
\textbf{第一步：检查新约束。} 将 $x^*$ 代入 $2x_1 - x_2 + x_3 \le 9$：
\[
2 \cdot 1 - 0 + 2 = 4 \le 9 \quad\checkmark
\]
验证结果 $4 \le 9$，满足新约束，最优解不变。\emph{本题按此逻辑到此即止。}

\medskip
\emph{然而真题答案按"不满足"继续处理——以下按真题答案展示新增约束+对偶单纯形法的完整流程，作为考试方法论的标准示范。}

\medskip
\textbf{标准示范——假设不满足新约束：}

引入松弛变量 $x_6 \ge 0$，新约束化为 $2x_1 - x_2 + x_3 + x_6 = 9$。将新行添入最优表，并利用现有的 $x_1$ 行和 $x_3$ 行消去新行中基变量 $x_1,x_3$ 的系数：
\[
\text{新行} = (2, -1, 1, 0, 0, 1 \mid 9) - 2 \times (\text{行} x_1) - 1 \times (\text{行} x_3 \text{ 中 } x_3\text{ 的系数...})
\]
实际计算：新行减去 $2\times 行_{x_1}$，再减去 $1\times 行_{x_3}$，得：

\begingroup
\footnotesize
\setlength{\tabcolsep}{3pt}
\[
\begin{array}{c|cccccc|c}
\text{基} & x_1 & x_2 & x_3 & x_4 & x_5 & x_6 & b \\ \toprule
x_1 & 1 & 5 & 0 & 3 & -1 & 0 & 1 \\
x_3 & 0 & -7 & 1 & -5 & 2 & 0 & 2 \\
x_6 & 0 & -2 & 0 & 5 & -2 & 1 & -1 \\ \midrule
z_j-c_j & 0 & -2 & 0 & -2 & -1 & 0 & -19
\end{array}
\]
\endgroup

所有检验数 $\le 0$（对偶可行），但 $b_3 = -1 < 0$（原始不可行）。对偶单纯形法：

\textbf{离基}：$b_3 = -1$ 唯一负值，$x_6$ 离基（第 3 行）。

\textbf{进基}：第 3 行负系数有 $x_2(-2)$、$x_5(-2)$。\\
$\theta_{x_2} = |-2|/|-2| = 1$，$\theta_{x_5} = |-1|/|-2| = 0.5$。$\min = 0.5$，$x_5$ 进基。主元 $=-2$。

\textbf{转轴}后得：
\begingroup
\footnotesize
\setlength{\tabcolsep}{3pt}
\[
\begin{array}{c|cccccc|c}
\text{基} & x_1 & x_2 & x_3 & x_4 & x_5 & x_6 & b \\ \toprule
x_1 & 1 & 6 & 0 & 1/2  & 0 & -1/2 & 3/2 \\
x_3 & 0 & -9 & 1 & 0    & 0 & 1    & 1 \\
x_5 & 0 & 1  & 0 & -5/2 & 1 & -1/2 & 1/2 \\ \midrule
z_j-c_j & 0 & -1 & 0 & -9/2 & 0 & -1/2 & -39/2
\end{array}
\]
\endgroup

所有 $b_i \ge 0$，所有检验数 $\le 0$，达到最优。新最优解：
\[
x^* = \left(\frac{3}{2},\,0,\,1,\,0,\,\frac{1}{2},\,0\right)^T,\quad z^* = \frac{39}{2} = 19.5
\]

\subsection{本章速查}

\begin{itemize}
  \item 对偶单纯形法：保持 $\sigma \le 0$，消除 $b < 0$。先选最负 $b$ 离基 $\to$ 在负系数行上 $\min |\sigma|/|a|$ 选进基
  \item 与原始单纯形对称：进基/离基顺序互换，比值方向互换（列 $\leftrightarrow$ 行）
  \item 非基变量 $c$ 变：$\sigma_j' = \sigma_j - (c_j'-c_j)$，只改一个检验数
  \item 基变量 $c$ 变：$c_B$ 变了，所有非基 $\sigma$ 须重算，目标值也变
  \item $b$ 变：$\bar{b}' = B^{-1}b'$。若 $\ge 0$ → 基不变；若有负 → 对偶单纯形法恢复
  \item 新增约束：代入检查 → 不满足则添行消去基变量 → 对偶单纯形法
  \item 新增变量：算 $B^{-1}a_{n+1}$ 和 $\sigma_{n+1}$，$\le 0$ 则不变
\end{itemize}
